\section{Conclusion}\label{sec:conclusion}
The four retrospective manuscript questions have conditional, evidence-based
answers. For Q1, aggregate ANWG orderings are positively associated across
all 18 source-pair/region comparisons ($\tau_b=0.551$--$1.000$), with
stronger Azure--Bailian agreement than BurstGPT comparisons. Absolute
utility and tied top sets differ markedly across sources, so ordinal
agreement alone does not establish equivalent service quality.

For Q2, 36 of 990 records pass the historical 10\% point-margin and nominal
interval rule, spanning five policy pairs. The endpoints reproduce, but
calibrated significance and FDR control are not established; alternative
family bookkeeping and grouped-window resampling each retain 30 under
their respective rules. The finding is a minority of observed practical
ordering changes in a specified simulator, not 36 confirmed physical
reversals.

For Q3, the historical median cross-metric tau is 0.419. Deterministic
ID agreement is 31.9\%, whereas maximizing sets intersect in 57.8\% of
conditions. Matching observed windows leaves the broad descriptive pattern
intact but changes some estimates and their target population. Saved SLO
sensitivities likewise constrain any objective-independent selection claim.

For Q4, most temporal comparisons are stable, with a BurstGPT LOW exception
at $\tau_b=-0.10$. Recovery depends on source and target: the 0.9 exact-order
criterion is reached at 30 windows for Azure and only at the trivial
40-window reference for the other sources. This is finite-reference
stability, not an independent-window requirement for future workloads.

The contribution is a traceable comparison design and a critical empirical
case study, not a new rank statistic or a universal scheduler recommendation.
Original synthetic-to-real and faithful native-ranking questions remain
unanswered. Policy-dependent service semantics and legacy physical-runner
defects prevent transporting these results to production serving without
separate comparability validation and new, prospectively designed evidence.
